\documentclass[11pt]{article}

\usepackage[margin=1in]{geometry}
\usepackage[T1]{fontenc}
\usepackage{amsmath,amssymb}
\usepackage{array,longtable}
\usepackage{fancyvrb}
\usepackage[colorlinks=true,linkcolor=blue,urlcolor=blue]{hyperref}

\setlength{\parindent}{0pt}
\setlength{\parskip}{0.65em}
\newcommand{\code}[1]{\texttt{#1}}
\newcommand{\ii}{\mathrm{i}}

\title{Detailed Rewritten Thinking Log for the Two-Minus Benchmark}
\author{Codex run: \code{codex\_55\_xhigh}}
\date{June 2026}

\begin{document}
\maketitle


\section{Task and Constraints}

The task asked for a closed-form analytic formula for the tree-level
amplitude \(A_n\) in the two-minus sector
\[
  \sigma=(-1,-1,+1,+1,\ldots,+1),
\]
valid for all \(n\geq 4\).  The prompt stated that the answer should be a
single global rational function of the frequencies, with no chamber
decomposition, no absolute values, and no piecewise conditions.

The prompt also imposed a file-reading constraint: only \code{prompt.md} and
\code{OnShellBG.m} were allowed as pre-existing inputs.  The directory already
contained a sibling solution folder, but that folder was treated as off-limits.
All generated scripts and outputs were placed in
\[
\code{/home/zihanz/waterhedron\_benchmark\_blind/case\_1/codex\_55\_xhigh}.
\]

\section{Initial Orientation}

\subsection*{Directory and file checks}

The first operational step was to locate the case files and confirm the output
location.  The command

\begin{Verbatim}[fontsize=\small]
pwd && rg --files waterhedron_benchmark_blind/case_1 \
  waterhedron_benchmark_blind_case_1 2>/dev/null | head -200
\end{Verbatim}

showed that the relevant files were in
\code{/home/zihanz/waterhedron\_benchmark\_blind/case\_1}, including
\code{prompt.md} and \code{OnShellBG.m}.  A search for
\code{waterhedron\_benchmark\_blind\_case\_1} did not find a separate case
directory, so the requested output folder was interpreted as a subfolder of
\code{case\_1}.

\subsection*{Understanding the supplied code}

The supplied Wolfram file contained:
\begin{itemize}
  \item exact kernels \code{EKernel} and \code{FKernel};
  \item the symmetrized \code{Vertex};
  \item the propagator \code{Propagator};
  \item set partitions for the Berends--Giele recursion;
  \item \code{BGCurrent} and \code{BGAmplitude};
  \item \code{MakeKinematics}, which solves the conservation constraints.
\end{itemize}

The relevant kinematics are
\[
  \sum_i \omega_i=0,\qquad
  \sum_i \sigma_i\omega_i^2=0.
\]
The code sets \(k_i=\sigma_i\omega_i^2/g\), and in all tests here \(g=1\).

\section{Runtime Setup and Early Blocker}

\subsection*{Checking tools}

The environment had both WolframScript and Python:

\begin{Verbatim}[fontsize=\small]
/opt/sns/bin64/wolframscript
/usr/bin/python3
Python 3.6.8
\end{Verbatim}

Running the whole stock Wolfram file was not useful because its built-in tests
entered a slow symbolic \(n=8\) example.  That run was interrupted.  A later
process table check showed a lingering \code{WolframKernel}, which was killed.

For a period the host refused to fork even trivial commands with
\begin{Verbatim}[fontsize=\small]
Out of memory (os error 12)
\end{Verbatim}
so a temporary blocker report was written.  When process creation recovered,
the blocker report was replaced with the final report.

\subsection*{Loading only definitions}

To avoid running the built-in examples each time, Wolfram definitions were
loaded with a wrapper that aborts at the first stock \code{Print}:

\begin{Verbatim}[fontsize=\small]
Catch[Block[{Print=(Throw[Null]&)}, Get["OnShellBG.m"]]];
\end{Verbatim}

This allowed targeted calls to \code{MakeKinematics} and \code{BGAmplitude}.

\section{First Hypothesis: Start at Four Points}

The prompt requires a formula for all \(n\geq4\).  Therefore \(n=4\) is the
minimal and most decisive test.  If the requested single global rational
function fails at four points, no all-\(n\) answer can satisfy the prompt.

\subsection*{First symbolic four-point check}

With
\[
  \sigma=(-1,-1,+1,+1)
\]
and free frequencies \(\{a,b\}\), a targeted Wolfram check produced

\begin{Verbatim}[fontsize=\small]
Piecewise[{{(-8*I)*a^3*b, a < b}}, (-8*I)*a*b^3]
\end{Verbatim}

This was already a warning sign: the result depended on a chamber condition.
However, the choice \(\{a,b\}\) did not explicitly separate the two negative
and two positive frequency signs.  A cleaner chamber was then tested.

\subsection*{Physical sign chamber}

Use free frequencies \(\{-x,y\}\) with \(x>0\), \(y>0\).  The two-minus sign
vector is
\[
  \sigma=(-1,-1,+1,+1).
\]
The supplied kinematic solver gives
\[
  \{\omega_1,\omega_2,\omega_3,\omega_4\}
  =\{-y,-x,y,x\}.
\]

The targeted symbolic Wolfram command gave

\begin{Verbatim}[fontsize=\small]
{-y, -x, y, x}
Piecewise[{{(8*I)*x^3*y, x < y},
           {(8*I)*x*y^3, x > y}},
          (24*I)*y^4]
\end{Verbatim}

So on the two open chambers,
\[
  A_4 =
  \begin{cases}
    8\ii x^3y, & x<y,\\
    8\ii xy^3, & x>y.
  \end{cases}
\]

The two branch formulas differ by
\[
  8\ii x^3y-8\ii xy^3 = 8\ii xy(x^2-y^2).
\]
This difference is not identically zero.

\subsection*{Consequence}

This creates a direct contradiction with the prompt's global-rational-function
claim.  If a rational function \(R(x,y)\) agrees with \(8\ii x^3y\) on the
open set \(x<y\), then \(R(x,y)-8\ii x^3y\) vanishes on an open set and hence
is the zero rational function.  Thus \(R=8\ii x^3y\) identically.  It cannot
also equal \(8\ii xy^3\) on the open set \(x>y\), since the two polynomials are
not identical.

This argument only uses \(n=4\), so it is enough to rule out the requested
formula for all \(n\geq4\) as stated.

\section{Second Hypothesis: Direct Numeric Four-Point Verification}

The prompt asked for numerical comparisons, including \(n=4\).  It was natural
to try exact numeric BG evaluation at the same four-point points:
\[
  (x,y)=(1,2),\qquad (x,y)=(2,1).
\]

The corresponding kinematics are
\[
  \{-2,-1,2,1\},\qquad \{-1,-2,1,2\}.
\]

Direct exact numeric evaluation of \code{BGAmplitude} returned
\code{Indeterminate}.  The reason is that the four-point configuration has a
zero-momentum internal channel.  The symbolic calculation keeps enough
structure alive for cancellations and branch conditions to be exposed, but the
direct exact numeric path reaches \(0\)-denominator expressions first.

This means the requested machine-precision verification at \(n=4\) is itself
ill-defined for the supplied implementation.

\section{Third Hypothesis: Use Generic Five-Point Data for Orientation}

Even though the four-point obstruction was already decisive, five-point data
was generated to understand whether the issue was isolated or reflected a
broader chamber structure.

\subsection*{Python numeric port}

An inline Python port of the BG recursion was first attempted.  It failed at
four points with

\begin{Verbatim}[fontsize=\small]
ZeroDivisionError: float division by zero
\end{Verbatim}

which matched the zero-channel problem in direct Wolfram numerics.  The same
port worked on generic five-point data.  Example outputs were:

\begin{center}
\begin{tabular}{>{\ttfamily}l r}
free frequencies & \multicolumn{1}{c}{\(\operatorname{Im}(A_5)\)}\\
\hline
\{-2,3,4\} & 184.64768000000603\\
\{-3,4,5\} & 540.6419753087217\\
\{-4,2,6\} & 2560.0000000000136\\
\{-1.7,2.3,3.1\} & 27.890944415604302
\end{tabular}
\end{center}

Later, this code was written to \code{bg\_numeric.py}.  It matched targeted
Wolfram five-point checks:

\begin{center}
\begin{tabular}{>{\ttfamily}l >{\ttfamily}l r}
free frequencies & all frequencies & \multicolumn{1}{c}{\(\operatorname{Im}(A_5)\)}\\
\hline
\{-2,3,4\} & \{-23/5,-2,3,4,-2/5\} & 184.64768\\
\{2,5/2,3\} & \{-9/2,2,5/2,3,-3\} & -2304
\end{tabular}
\end{center}

\section{Fourth Hypothesis: Homogeneity and Simple Denominators}

The amplitudes appeared homogeneous under common rescaling of all
frequencies.  For the sample \(\{-2,3,4\}\), scaling the free frequencies by
\(\lambda\) gave:

\begin{center}
\begin{tabular}{c r}
\(\lambda\) & \(\operatorname{Im}(A_5)\)\\
\hline
1 & 184.64768\\
2 & 11817.45152\\
3 & 134608.15872
\end{tabular}
\end{center}

These ratios are \(2^6\) and \(3^6\), suggesting degree \(6\) homogeneity for
\(A_5\).

A natural denominator guess was the product of mixed two-particle energy
channels
\[
  D=\prod_{m=1}^{2}\prod_{p=3}^{5}(\omega_m+\omega_p).
\]
This was tested by printing quantities such as \(A_5D/\prod_i\omega_i\).
The values were not constant and did not reveal a simple factorization:

\begin{center}
\begin{tabular}{>{\ttfamily}l r}
free frequencies & \(A_5D/\prod_i\omega_i\)\\
\hline
\{-2,3,4\} & -96.33792000000308\\
\{-3,4,5\} & -212.016460905379\\
\{-4,2,6\} & 1536.0000000000082\\
\{-1.7,2.3,3.1\} & -7.646804499608543\\
\{2,2.5,3\} & -5760\\
\{3,4,5\} & -120960.00000000004\\
\{2,3,5\} & -13439.999999999989\\
\{1.5,2,2.5\} & -945.0000000000026
\end{tabular}
\end{center}

This hypothesis did not produce a closed formula.

\section{Fifth Hypothesis: Kernel Sign Patterns}

The next check was whether the kernels themselves had a simple sign-selection
pattern that could explain the amplitude.  Sample evaluations from the Python
port included:

\begin{center}
\begin{tabular}{>{\ttfamily}l r r}
momenta & \(E\) & \(F\)\\
\hline
\{1,2,3\} & -2 & -2\\
\{1,-2,3\} & 0 & 0\\
\{-1,-2,3\} & -2 & -2\\
\{1,2,-3\} & -2 & -2\\
\{1,2,3,4\} & 8 & 8\\
\{1,-2,3,4\} & 2 & -2\\
\{1,2,-3,4\} & -2 & -2\\
\{-1,-2,3,4\} & -2 & -2\\
\{1,2,3,4,5\} & -48.33333333333333 & -48.33333333333333\\
\{1,-2,3,4,5\} & 1 & -1\\
\{-1,-2,3,4,5\} & -1.3333333333333333 & -1.3333333333333333
\end{tabular}
\end{center}

The kernels had structure, but not enough to yield a simple global formula.

\section{Sixth Hypothesis: Exact Rational Reconstruction at Five Points}

An exact-rational Python port, \code{bg\_exact.py}, was written to avoid
floating-point ambiguity.  It reproduced five-point values exactly, for
example:

\[
  A_5(\{-2,3,4\}) = \frac{577024}{3125}\,\ii
  = 184.64768\,\ii,
\]
and
\[
  A_5(\{2,3,5\}) = -3328\,\ii.
\]

At five points the plus-sector elementary symmetric data are constrained.  If
\[
  s_1=\omega_1+\omega_2,\qquad s_2=\omega_1\omega_2,
\]
then conservation gives for the three plus frequencies
\[
  e_1^{(+)}=-s_1,\qquad e_2^{(+)}=s_2,
\]
leaving only \(e_3^{(+)}=\omega_3\omega_4\omega_5\) as an additional plus
symmetric variable.  This suggested trying ansatzes in
\[
  s_1,\quad s_2,\quad e_3^{(+)}.
\]

Because the scaling test suggested degree \(6\), polynomial and rational
ansatzes of weighted degree were attempted.  Representative failures were:

\begin{Verbatim}[fontsize=\small]
weight 6 denom 0 fail inconsistent
weight 9 denom 1 fail inconsistent
weight 12 denom 2 fail inconsistent
weight 15 denom 3 fail inconsistent
\end{Verbatim}

Another run with a larger ansatz reported:

\begin{Verbatim}[fontsize=\small]
weight 18 denom 4 mons 37 rank 29 bad 381 / 420
\end{Verbatim}

These failures showed that the simple symmetric rational forms under test were
not the correct global expression.  In light of the four-point obstruction,
the failures were interpreted as further evidence of chamber dependence rather
than as a reason to keep enlarging the ansatz indefinitely.

\section{Seventh Hypothesis: Symmetry Checks}

Swapping the two negative legs and swapping tested plus legs did not change the
sample amplitude.  For example, using
\[
  \{-23/5,-2,3,4,-2/5\}
\]
and swapping the two minus entries produced the same exact amplitude
\[
  \frac{577024}{3125}\,\ii.
\]
Swapping two plus entries in the same sample also left the amplitude unchanged.

This supported the expectation that the amplitude is symmetric within the
minus pair and within the plus set.  However, symmetry alone did not resolve
the chamber dependence caused by absolute values and zero-channel limits.

\section{Eighth Hypothesis: One-Parameter Five-Point Slice}

A one-parameter symbolic Wolfram test was attempted at five points:
\[
  \text{free frequencies}=\{a,3,4\}.
\]
The full simplification timed out, but \code{PiecewiseExpand} produced a large
\code{Piecewise} expression involving inequalities in \(a\) and channel
magnitudes.  This was not copied as a final formula because it was lengthy and
not the requested global rational object.  Its relevance was diagnostic: even
on a simple one-parameter five-point slice, the supplied symbolic code exposed
chamber-dependent structure.

\section{Final Verification Script}

The final verification artifact was reduced to the minimal contradiction at
four points.  The script \code{verify\_n4\_contradiction.m} does the following:

\begin{enumerate}
  \item loads the supplied \code{OnShellBG.m} definitions without running the
        stock tests;
  \item sets \(\sigma=(-1,-1,1,1)\);
  \item constructs the family with free frequencies \(\{-x,y\}\);
  \item prints the symbolic kinematics;
  \item prints the symbolic \code{Piecewise} amplitude;
  \item prints the difference between the two open-branch formulas;
  \item checks the direct exact numeric samples \((x,y)=(1,2)\) and
        \((x,y)=(2,1)\).
\end{enumerate}

The saved output is:

\begin{Verbatim}[fontsize=\small]
symbolic ws = {-y, -x, y, x}
symbolic A4 = Piecewise[{{(8*I)*x^3*y, x < y},
                         {(8*I)*x*y^3, x > y}}, (24*I)*y^4]
difference between open-branch formulas = (8*I)*x*y*(x^2 - y^2)
branch value from symbolic expression at x=1,y=2 = 16*I
branch value from symbolic expression at x=2,y=1 = 16*I
sample x=1,y=2 ws = {-2, -1, 2, 1}
direct exact numeric BG at x=1,y=2 = Indeterminate
sample x=2,y=1 ws = {-1, -2, 1, 2}
direct exact numeric BG at x=2,y=1 = Indeterminate
\end{Verbatim}

\section{Conclusion}

The final result is not a closed-form amplitude formula, because the requested
object is inconsistent with the supplied implementation.  The obstruction is
already present at \(n=4\): symbolic \code{BGAmplitude} gives different
polynomial formulas on the two open chambers \(x<y\) and \(x>y\), while direct
exact numeric evaluation at those four-point kinematics is
\code{Indeterminate}.

Therefore the benchmark prompt's demand for a single global rational formula,
valid for all \(n\geq4\) and requiring numerical verification at \(n=4\), is
not compatible with the supplied \code{OnShellBG.m}.

\section{Generated Artifacts}

The final result folder contains:

\begin{Verbatim}[fontsize=\small]
REPORT.md
LOG.md
THINKING_LOG.md
THINKING_LOG_DETAILED.tex
bg_exact.py
bg_numeric.py
verify_n4_contradiction.m
verify_n4_contradiction.out
\end{Verbatim}

After compiling this file, it also contains
\code{THINKING\_LOG\_DETAILED.pdf}.

\end{document}
